% Source: source/普通高考/2008/2008新课标文(海南,宁夏).pdf, page 1, question 6.
% pdfimages -list found no embedded bitmap; the program flowchart is vector artwork.
% Node -> directed-edge list: 开始 -> 输入 a,b,c -> x=a -> b>x?;
% 否 -> second (blank) decision -> 输出 x -> 结束;
% 是 -> x=b -> left return arrow -> shared stem above the second decision;
% second 是 -> x=c -> left return arrow -> shared stem above 输出 x.
% All connectors are solid; branch labels are attached to their corresponding edges.

\begin{tikzpicture}[
    >=Stealth,
    x=1cm,
    y=0.86cm,
    line width=0.45pt,
    line cap=round,
    line join=round,
    font=\normalsize,
    every node/.style={anchor=center, align=center,
        execute at begin node={\setlength{\parindent}{0pt}}},
    flowbox/.style={draw, anchor=center, align=center, inner xsep=3pt,
        inner ysep=2pt, execute at begin node={\setlength{\parindent}{0pt}}},
    terminal/.style={flowbox, rounded corners=0.16cm,
        minimum width=1.35cm, minimum height=0.68cm},
    process/.style={flowbox, minimum height=0.68cm},
    decision/.style={flowbox, diamond, aspect=3.4, inner sep=0pt,
        minimum width=3.45cm, minimum height=1.02cm},
    io/.style={flowbox, trapezium, trapezium left angle=72,
        trapezium right angle=108, minimum height=0.76cm},
    branchlabel/.style={anchor=center, align=center, inner sep=0pt,
        execute at begin node={\setlength{\parindent}{0pt}}},
]
    \path[use as bounding box] (-2.95,-1.05) rectangle (4.20,11.45);

    \node[terminal] (start) at (0,10.85) {开始};
    \node[io, minimum width=3.55cm] (input) at (0,9.40) {输入 $a,b,c$};
    \node[process, minimum width=1.65cm] (seta) at (0,7.95) {$x=a$};
    \node[decision] (first) at (0,6.42) {$b>x$};
    \node[process, minimum width=1.75cm] (setb) at (3.30,5.00) {$x=b$};
    \node[decision] (second) at (0,3.60) {};
    \node[process, minimum width=1.75cm] (setc) at (3.30,2.58) {$x=c$};
    \node[io, minimum width=2.15cm] (output) at (0,0.90) {输出 $x$};
    \node[terminal] (finish) at (0,-0.50) {结束};

    % Main chain and the first decision branch.
    \draw[->] (start.south) -- (input.north);
    \draw[->] (input.south) -- (seta.north);
    \draw[->] (seta.south) -- (first.north);
    \draw[->] (first.south) -- node[branchlabel,left=2pt] {否} (second.north);
    \draw[->] (first.east) -- node[branchlabel,above=1pt] {是} (3.30,6.42)
        -- (setb.north);
    % Return path merges into the shared stem above the second decision.
    \draw[->] (setb.south) -- (3.30,4.40) -- (0,4.40);
    \draw (0,4.40) -- (second.north);

    % Second decision branch and the output path.
    \draw[->] (second.south) -- node[branchlabel,left=2pt] {否} (output.north);
    \draw[->] (second.east) -- node[branchlabel,above=1pt] {是} (3.30,3.60)
        -- (setc.north);
    % Return path merges into the shared stem above the output node.
    \draw[->] (setc.south) -- (3.30,1.95) -- (0,1.95);
    \draw (0,1.95) -- (output.north);
    \draw[->] (output.south) -- (finish.north);
\end{tikzpicture}
